\lecture{6}
\title{Motion Planning}

\begin{document}

\subsection{Motion Planning: Setup}

\textbf{Definition.}
Some reminders of some definitions:
\begin{itemize}
    \item The \textbf{work space} (or \textbf{Cartesian space}) of a robot
    is the set of points reachable by an end-effector.
    \begin{itemize}
        \item For example, one might use $(x, y, \theta)$ to parametrize
        the work space of a two-rod 2D Kuka IIWA.
    \end{itemize}
    \item The \textbf{configuration space} of a robot
    is the set of reachable configurations of the robot---not just
    the end-effector!
    \begin{itemize}
        \item A configuration is denoted $\mathbf{q} \in C$,
        where $C$ is the set of all possible configurations.
        \item For example, one might use $(q_1, q_2)$ to parametrize
        the configuration space of a two-rod 2D Kuka IIWA.
        \begin{itemize}
            \item The topology of the configuration space
            might not be Cartesian; in this example,
            it's $S^1 \times S^1$ (a torus).
        \end{itemize}
    \end{itemize}
\end{itemize}

\textbf{Definition.}
In motion planning, we assume that we have:
\begin{itemize}
    \item A world $W$.
    \item An obstacle region $O \subset W$.
    \item A robot $A$ defined in $W$ (i.e., $A \subset W$).
    \item We seek a collision-free path from
    $\mathbf{q}_I \in C$ to $\mathbf{q}_G \in C$.
\end{itemize}

\subsection{Motion Planning: Obstacle Space}

Given a robot $A$ and an object $O$,
how can we determine the set of positions of $A$
in which $A$ and $O$ do not intersect?

\textbf{Definition.}
We define the \textbf{obstacle space} $C_{\mathrm{obs}} \subset C$ by
$C_{\mathrm{obs}} := \{q \in C \mid A(q) \cap O \neq \emptyset\}$,
in addition to $C_{\mathrm{free}} = C \setminus C_{\mathrm{obs}}$.

\begin{remark}
    For some reason, we care about whether $C_{\mathrm{free}}$ is closed
    in the topology of $C$.  If not, we must use
    the closure of $C_{\mathrm{free}}$.
\end{remark}

In simple translational cases,
we can write $C_{\mathrm{obs}} = O \ominus A$,
where $\ominus$ is the
\href{https://en.wikipedia.org/wiki/Minkowski_addition}{Minkowski difference}.

\begin{center}
    \frame{\includegraphics[width=0.4\textwidth]{lecture-06/minkowski-difference.png}}
\end{center}

The Minkowski difference can be efficiently computed using the
\textbf{Star algorithm}---details omitted.

But for nontranslational cases, the obstacle space can be harder to compute.
For example, the robot might be allowed to rotate,
or the configuration space might be non-Cartesian.

\begin{center}
    \frame{\includegraphics[width=0.4\textwidth]{lecture-06/rotating.png}}
    \frame{\includegraphics[width=0.4\textwidth]{lecture-06/two-rod-obstacles.png}}
\end{center}

In some of these cases, the runtime for computing $C_{\mathrm{obs}}$
might be exponential!

The upshot of defining $C_{\mathrm{obs}}$ is that
we can now define the problem of motion planning
entirely within the configuration space $C$.

\textbf{Definition.}
In motion planning, we assume that we have:
\begin{itemize}
    \item A configuration space $C$, plus
    an obstacle space $C_{\mathrm{obs}} \subset C$
    and a free space $C_{\mathrm{free}} := C \setminus C_{\mathrm{obs}}$.
    \item We seek a continuous path $\tau \colon [0, 1] \to C_{\mathrm{free}}$
    with $\tau(0) = \mathbf{q}_I$ and $\tau(1) = \mathbf{q}_G$.
\end{itemize}

\subsection{Combinatorial Planning}

\textbf{Definition.}
A motion planner is \textbf{complete} if
it always finds a solution if one exists
and returns a failure in \emph{bounded} time if not.

The combinatorial algorithms we discuss in this section will be complete,
and they will all employ \textbf{road maps}.

\textbf{Definition.}
Given a configuration space $C$ and a free space $C_{\mathrm{free}}$,
a \textbf{road map} $R$ is a graph whose vertices are
configurations in $C_{\mathrm{free}}$ and whose edges are
collision-free paths between vertices.
This graph must be such that:
\begin{itemize}
    \item \textbf{(Accessibility)}
    For any starting point $q_I \in C_{\mathrm{free}}$,
    it is easy to find a collision-free path $q_I \to q'$
    for some $q' \in R$.
    \item \textbf{(Departability)}
    For any ending point $q_G \in C_{\mathrm{free}}$,
    it is easy to find a collision-free path $q' \to q_G$
    for some $q' \in R$.
    \item \textbf{(Connectivity)}
    If $q_1, q_2 \in R$ are connected in $C_{\mathrm{free}}$,
    then they are connected in the graph of $R$.
\end{itemize}
Here's an example of how a road map
might let us find a collision-free path $q_I \to q_G$.
\begin{center}
    \frame{\includegraphics[width=0.25\textwidth]{lecture-06/roadmap.png}}
\end{center}

Here are some different strategies for constructing road maps.

\begin{quote}
    \textbf{Visibility Graph.}
    Suppose the configuration space is $\mathbb{R}^2$,
    and every obstacle is a convex polygon.
    (If an obstacle is concave,
    you can work around this by breaking it up
    into convex pieces.)

    Then for any $q_I, q_G \in C_{\mathrm{free}}$,
    to find a path $q_I \to q_G$, we simply:
    \begin{itemize}
        \item Precompute for any two vertices of two polygons,
        whether the straight-line path between them is collision-free.
        \item Determine the set of all polygon vertices $q$
        for which $q_I \to q$ is collision-free
        (and separately for $q \to q_G$, too).
        \item Apply a shortest-path $O(n^2 \log n)$ algorithm (e.g., A*)
        to find the shortest collision-free path $q_I \to q_G$.
    \end{itemize}
    \begin{center}
        \frame{\includegraphics[width=0.25\textwidth]{lecture-06/visibility-graph.png}}
    \end{center}
    It turns out that this algorithm in fact
    always finds the \emph{shortest} collision-free path $q_I \to q_G$.
\end{quote}

\begin{remark}
    In computational geometry,
    people have found more efficient algorithms
    for computing visibility graphs.
    In practice, though,
    sampling-based motion planning algorithms are preferred---as
    we'll see later on.
\end{remark}

\begin{quote}
    \textbf{Voronoi Diagram.}
    Just construct the
    \href{https://en.wikipedia.org/wiki/Voronoi_diagram}{Voronoi Diagram}
    of the obstacles.
    \begin{center}
        \frame{\includegraphics[width=0.25\textwidth]{lecture-06/voronoi-diagram.png}}
    \end{center}
    The roadmap graph is then exactly the ``boundaries''
    of the Voronoi diagram.
    In a sense, this lets you keep a maximum distance away
    from any obstacle throughout your collision-free path.
\end{quote}

One last algorithm archetype.

\begin{quote}
    \textbf{Cell Decomposition.}
    Split the configuration space into cells,
    and then find paths between the cells.

    For example, here's a cell decomposition by
    $x$-coordinate of the obstacles.
    \begin{center}
        \frame{\includegraphics[width=0.4\textwidth]{lecture-06/x-coordinate-cells.png}}
    \end{center}
    There are also \emph{approximate} cell decomposition approaches;
    for example, here's a quadtree decomposition approach.
    \begin{center}
        \frame{\includegraphics[width=0.4\textwidth]{lecture-06/quadrtree-decomposition.png}}
    \end{center}
    The idea is that cells are either empty, full, or mixed,
    and that ``mixed'' cells can be zoomed into as needed.
\end{quote}

As mentioned earlier, these are all \textbf{complete} algorithms
(approximate cell decompositions only up to resolution).
But their disadvantages are that they become intractable
as the obstacle complexity or dimensionality of the $C$-space increases.
(And this happens very often in robotic manipulation!)

\subsection{Sampling-Based Planning}

\textbf{Definition.}
A (randomized) motion planner is \textbf{probabilistically complete}
if, assuming a solution exists,
the probability of finding that solution in $t$ time
tends to one as $t \to \infty$.

The sampling-based algorithms we discuss in this section
will be probabilistically complete.  There are two kinds:
\begin{itemize}
    \item \textbf{Multi-query:}
    The algorithm performs some precomputation---assumes
    obstacles are fixed across queries.
    \item \textbf{Single-query:}
    The algorithm does not perform precomputation---useful
    when obstacles move over time.
\end{itemize}

Here's the most common multi-query algorithm.

\begin{quote}
    \textbf{Probabilistic Roadmap (PRM).}
    Construct the roadmap $R = (V, E)$ procedurally by repeating the following algorithm.
    \begin{enumerate}
        \item Randomly pick a point $q \in C$
        and check if $q \in C_{\mathrm{free}}$.
        If it isn't, try again.
        \item If it is, add $q$ to $V$,
        and for each $q' \in V$ with $\|q - q'\| < d$,
        add the edge $q \leftrightarrow q'$ to $E$
        if that path is collision-free.
    \end{enumerate}
    The whole point of PRM is to avoid having to compute $C_{\mathrm{obs}}$.
    Instead, for each $q \in C$, it uses a custom collision-detector
    to check whether $q \in C_{\mathrm{free}}$.
    (Efficient collision-detectors exist,
    because \emph{video games} require them.)
    \begin{center}
        \includegraphics[width=0.3\textwidth]{lecture-06/PRM.png}
    \end{center}
    As shown above, though, the paths outputted by PRM
    are not very short.  Here are some improvements to PRM:
    \begin{itemize}
        \item Find \emph{shortcuts}---between two points on the path,
        ask if the straight line between them is collision-free.
        \begin{center}
            \frame{\includegraphics[width=0.3\textwidth]{lecture-06/shortcut.png}}
        \end{center}
        \item Use paths generated by PRM as starting paths
        for the optimization algorithms described in
        \href{/6.4210/lectures/5}{Lecture 5}.
    \end{itemize}
\end{quote}

Unfortunately, PRM may fail with nonzero probability
because the roadmaps it produces might not satisfy
\emph{Accessibility and Departability} or \emph{Connectivity},
as shown below.

\begin{center}
    \frame{\includegraphics[width=0.6\textwidth]{lecture-06/PRM-failures.png}}
\end{center}

But PRM just works very well in practice because---even in
high dimensions---collision-detection is still very fast.

Sometimes the obstacles do move around, though,
in which case we require single-query algorithms.

\begin{quote}
    \textbf{Rapidly-Exploring Random Trees (RRT).}
    We gradually add points to a connected tree
    in the following way.
    \begin{enumerate}
        \item Randomly pick a point $q_{\mathrm{rand}} \in C$
        and check if $q_{\mathrm{rand}} \in C_{\mathrm{free}}$.
        If it isn't, try again.
        \begin{itemize}
            \item With probability $p$ (the \emph{goal bias}),
            our random point $q_{\mathrm{rand}}$
            will just be $q_{\mathrm{rand}} = q_G$,
            the goal.
        \end{itemize}
        \item (\textsc{Extend})
        If it is, then do the following:
        \begin{itemize}
            \item Find the nearest point $q' \in T$ to $q_{\mathrm{rand}}$.
            \item Compute the point $q_{\mathrm{new}}$
            that is $\eta$ of the way between $q'$ and $q_{\mathrm{rand}}$,
            for some fixed constant $\eta \in (0, 1)$.
            \item If the segment $q' \to q_{\mathrm{new}}$ lies in $C_{\mathrm{free}}$,
            then append $q_{\mathrm{new}}$ to $T$
            and connect $q'$ and $q_{\mathrm{new}}$ with
            a collision-free path.
        \end{itemize}
    \end{enumerate}
    You can see this algorithm in action
    using the (warning: AI-generated) widget below.

    \includewidget{lecture-06/rrt}

    Note that this algorithm is \emph{not} the same
    as a plain random walk---a random walk
    doesn't tend to go anywhere,
    whereas RRT will eventually tend to fill the space.

    A natural speed-up is \textbf{RRT-Connect},
    in which we grow a tree from \emph{both} $q_I$ and $q_G$
    bidirectionally.
    \begin{enumerate}
        \item Grow trees $T_A$ and $T_B$ out of $q_I$ and $q_G$,
        respectively.
        \item Suppose RRT adds a new vertex $q_{\mathrm{new}}$ to $T_A$
        (\textsc{Extend}).  Then:
        \begin{itemize}
            \item Take the nearest point $q' \in T_B$ to $q_{\mathrm{new}}$
            and take steps (\textsc{Extend}) from $q'$ to $q_{\mathrm{new}}$
            until a collision occurs.
            \item Add any taken steps to the tree $T_B$.
            Then swap $T_A$ and $T_B$.
            \begin{itemize}
                \item Allegedly, taking $T_A$ to always be the smaller tree
                helps even more.
            \end{itemize}
        \end{itemize}
        \item Repeat until the steps taken from $q'$ to $q_{\mathrm{new}}$
        reach $q_{\mathrm{new}}$ without any collisions.
    \end{enumerate}
    You can see in the widget above that
    RRT-Connect is much better than RRT for
    navigating narrow passages.

    One last improvement is \textbf{RRT*},
    which favors \emph{shorter} paths between $q_I$ and $q_G$.
    Handwavily, it adapts RRT in two ways:
    \begin{itemize}
        \item Connect each $q_{\mathrm{new}}$ to whichever nearby $q' \in T$
        gives the shortest path from $q_I$---not necessarily
        the $q'$ closest to $q_{\mathrm{new}}$.
        \item Afterwards, \emph{rewire} $T$:
        if a nearby $q' \in T$
        would have a shorter path from $q_I$ through $q_{\mathrm{new}}$,
        reroute it through $q_{\mathrm{new}}$.
    \end{itemize}
    RRT* keeps running even after it reaches $q_G$
    so that even shorter paths can be found over time.
\end{quote}

One advantage of RRT is that it works even for \textbf{non-holonomic} robots.

\textbf{Definition.}
A robot is \textbf{holonomic} if
it can freely move in any direction in its configuration space.
\begin{itemize}
    \item The simplified 2D two-rod Kuka IIWA
    is a holonomic robot.
    \item A car is not a holonomic robot;
    it cannot translate its body in a lateral direction.
\end{itemize}

Combinatorial planning algorithms
ask whether you can move from $q_1$ to $q_2$
in the configuration space,
which is not easily answerable for
non-holonomic robots.

But RRT only asks whether, starting from $q'$,
you can get \emph{closer} to $q_{\mathrm{rand}}$,
thereby landing on $q_{\mathrm{new}}$.
This can be answered for non-holonomic robots
by forward-simulating various inputs
and checking whether any results
land closer to $q_{\mathrm{rand}}$.

Spoilers: it turns out that RRT (or variants of it)
just perform the best in practice.

\end{document}